\section{Cross-notation equivalence and conversion}\label{sec:conversion-framework}

The forward maps of Sections~\ref{sec:petri-nets}--\ref{sec:bpmn} are model-to-model
transformations that carry each notation into a common framework. Equality of the
resulting signatures is then a model-equivalence check that can validate such a
transformation, confirming that two models agree as processes.

Sections~\ref{sec:petri-nets}--\ref{sec:bpmn} establish one direction, each notation mapping
canonically to a signature $\Sigma$, and hence to $\mathbf{F}(\Sigma)$, by one engine at different instance maps
(Remark~\ref{rem:ptree-recovers-pn}). Two questions follow, when two models, possibly in
different notations, produce equal signatures and what that equality certifies, and whether a
model can be recovered from a signature. We settle the first in
Section~\ref{subsec:equiv-certificate}. The second, and any full round-trip faithfulness, is
developed only as far as stating what each reconstruction requires
(Sections~\ref{subsec:reverse-map}--\ref{subsec:open-problems}) and is left as future work.

\begin{remark}[Notation-independence of the construction]
  \label{rem:ptree-recovers-pn}
  The notation constructions of Sections~\ref{sec:petri-nets}--\ref{sec:bpmn} are one engine at
  different instance maps. Each instance map reads typed wires $(x,w,t)$, treats silent leaves or
  transitions as transparent mediators with no elimination step, and composes along typed wires. A
  given object-centric behaviour therefore receives the same minimal signature whether presented as
  an OCPN, an OC causal net, or an OC process tree. The presentations differ only in how contexts
  are derived, from Petri place flow, from causal binding sets, or from operator contexts filtered
  by $D_{PT}$. A process tree also reaches its signature indirectly, by van Detten's OCPT-to-OCPN
  translation~\cite{vandettenDiscoveringCompactLive2024} composed with the OCPN construction.
  Provided that translation preserves behaviour, the indirect route and the direct construction of
  Section~\ref{sec:process_trees} yield the same signature, a commuting triangle that reads the
  notation-independence as a consistency check.
\end{remark}

\subsection{Signature equality as an equivalence certificate}\label{subsec:equiv-certificate}

Fix a process-mining notation $\mathcal{N}$ with the forward compilation $m\mapsto\Sigma(m)$ of
Definition~\ref{def:generator-cospan-general}, which consults no modelling convention. Write $C(m)$ for
the \emph{context structure} of $m$. The context structure is the family, indexed by activities, of the contexts
$c=(P,S)\in\mathrm{Contexts}(a)$ that the instance map reads off $m$'s compiled AND/OR graph
(Definition~\ref{def:contexts-general}), each carrying its constraint system $\Lambda_{a,c}$
(Definition~\ref{def:multiplicity-decoration}), compared by its solutions. All
four instance maps emit contexts in this one format, over the shared activity set and object types,
so $C(m)$ is defined in a notation-independent space (Remark~\ref{rem:ptree-recovers-pn}).

Fix a selection $\mathcal{S}$ in the sense of Section~\ref{subsec:theorem}, by default the
connected process diagrams, and call $L_{\mathcal{S}}(\Sigma)$, the union of the trace sets
$\Gamma(f)$ of Definition~\ref{def:trace-map} over $f\in\mathcal{S}(\Sigma)$, the trace language
of the model.

\begin{definition}[Structural and trace equivalence]\label{def:str-tr-equiv}
  Two models $m_1,m_2$, possibly in different notations, are \emph{structurally equivalent},
  $m_1\equiv_{\mathrm{str}}m_2$, if they induce the same context structure, $C(m_1)=C(m_2)$. They
  are \emph{trace-equivalent}, $m_1\equiv_{\mathrm{tr}}m_2$, if their trace languages coincide,
  $L_{\mathcal{S}}(\Sigma(m_1))=L_{\mathcal{S}}(\Sigma(m_2))$.
\end{definition}

\begin{lemma}[The signature is a complete invariant of the context structure]\label{lem:sig-complete}
  For models in any of the notations of Sections~\ref{sec:petri-nets}--\ref{sec:bpmn},
  \[
    \Sigma(m_1)=\Sigma(m_2) \quad\Longleftrightarrow\quad m_1\equiv_{\mathrm{str}}m_2,
  \]
  up to canonical generator relabelling. Structural equivalence is therefore decidable by direct
  comparison of the finite generator sets, across notations.
\end{lemma}

\begin{proof}
  The generate step (Definition~\ref{def:generator-cospan-general}) sends each context
  $c=(P,S)\in\mathrm{Contexts}(a)$ to the generator $g_{a,c}$ with label $\ell(a)$, input legs the
  typed wires $(x,w,a)\in P$, output legs the typed wires $(a,w,y)\in S$, and constraint system
  $\Lambda_{a,c}$. Because every typed wire names both endpoints and its set of types
  (Definition~\ref{def:contexts-general}), distinct contexts give distinct generators, so this map is injective. Reading each generator's boundary
  recovers its context, so it is a bijection between context structures and canonically-labelled
  signatures. Equal signatures therefore correspond exactly to equal context structures.
\end{proof}

\begin{lemma}[Soundness]\label{lem:equiv-soundness}
  $m_1 \equiv_{\mathrm{str}} m_2$ implies $m_1 \equiv_{\mathrm{tr}} m_2$.
\end{lemma}

\begin{proof}
  By Lemma~\ref{lem:sig-complete}, $\equiv_{\mathrm{str}}$ gives $\Sigma(m_1)=\Sigma(m_2)$, and Theorem~\ref{thm:canonical-presentation} gives
  $L_{\mathcal{S}}(\Sigma(m_1))=L_{\mathcal{S}}(\Sigma(m_2))$.
\end{proof}

\begin{lemma}[Strictness]\label{lem:equiv-strictness}
  $\equiv_{\mathrm{str}}$ is strictly finer than $\equiv_{\mathrm{tr}}$. Trace equivalence does not
  imply structural equivalence.
\end{lemma}

\begin{proof}
  Take one object type and weight $1$ throughout. Let $m_1$ be the sequence $a$ then $b$, and let
  $m_2$ offer three exclusive alternatives, $a$ then $b$, $a$ alone and $b$ alone, its occurrences
  of each activity being distinct nodes with the same label. Writing $(x,y)$ for the typed wire
  $(x,\{\ast\},y)$, $\Sigma(m_1)$ has the generators $a:(\bot,a)\to(a,b)$ and $b:(a,b)\to(b,\top)$.
  $\Sigma(m_2)$ has these two and also $a:(\bot,a)\to(a,\top)$ and $b:(\bot,b)\to(b,\top)$. The
  connected process diagrams of both signatures give the language $\{a,b,ab\}$, so
  $m_1\equiv_{\mathrm{tr}}m_2$, yet $C(m_1)\ne C(m_2)$ and $m_1\not\equiv_{\mathrm{str}}m_2$. The
  context structure records which occurrences can be joined to which, and the language loses
  that information.
\end{proof}

\begin{theorem}[Cross-notation equivalence certificate]\label{thm:signature-equivalence}
  Let $\mathcal{N}$ be any of the notations of Sections~\ref{sec:petri-nets}--\ref{sec:bpmn}.
  Then three statements hold.
  (i)~$\Sigma(m)$ is determined by $m$ alone, silent $\tau$-mediators contributing no generator and
  multiplicities carried in the constraint systems $\Lambda_{a,c}$, so no elimination step or routing convention is
  consulted. (ii)~$\Sigma(m_1)=\Sigma(m_2)$ decides $m_1\equiv_{\mathrm{str}}m_2$, cross-notation, by
  finite generator comparison (Lemma~\ref{lem:sig-complete}). The comparison is equality
  of $\Sigma\setminus\Sigma_\gamma$. The start and end generators of $\Sigma_\gamma$ state an
  assumption of the analysis (Section~\ref{subsec:analysis}) and take no part, so two models written
  in notations that open and close a system differently are compared without either being
  rewritten. (iii)~$\equiv_{\mathrm{str}}$ certifies $\equiv_{\mathrm{tr}}$ for every selection
  $\mathcal{S}$, by Theorem~\ref{thm:canonical-presentation}, and the converse fails
  (Lemmas~\ref{lem:equiv-soundness} and~\ref{lem:equiv-strictness}). In particular, since $\Sigma$ is convention-free,
  $\Sigma(m_1)\ne\Sigma(m_2)$ witnesses a structural difference that no encoding choice
  can produce, even when $m_1\equiv_{\mathrm{tr}}m_2$.
\end{theorem}

\begin{proof}
  Point (i) is the instance-map classification of silent transitions as mediators, which the walk
  passes through, with multiplicities kept as the constraint systems $\Lambda_{a,c}$
  (Section~\ref{subsec:firing-graphs}, Definitions~\ref{def:contexts-general}
  and~\ref{def:multiplicity-decoration}). Point (ii) is
  Lemma~\ref{lem:sig-complete}, and point (iii) is Lemmas~\ref{lem:equiv-soundness}
  and~\ref{lem:equiv-strictness}, whose contrapositive gives the negative certificate.
\end{proof}

For a practitioner the certificate answers whether two discovered models have the same structure,
and equal signatures then certify that they have the same trace language. It is strictly finer
than trace equivalence, as the sequence against choice case of
Lemma~\ref{lem:equiv-strictness} shows, so a pair the trace languages identify may still be
separated on structure. Inclusion gives the ordered form of the same check. When
$\Sigma(m_1)\subseteq\Sigma(m_2)$, every trace $m_1$ gives under a monotone selection is a trace of
$m_2$ (Corollary~\ref{cor:signature-inclusion}), and the generators of $\Sigma(m_2)$ outside
$\Sigma(m_1)$ name where $m_2$ allows more. Section~\ref{sec:validation} finds exactly this for two
discovered models, whose signatures nest with the causal net's inside the Petri net's.

\paragraph{Granularity.}
Structural equivalence is equality of context structures, which is coarser than isomorphism of the
AND/OR graphs. An OR gateway and its XOR/AND expansion are identified, as are two silent-mediator arrangements
with the same activity neighbourhoods. $\Sigma$ records the behavioural space, not the
syntax. Read as a model-equivalence relation, this coarseness compares discovered models over
their behavioural space rather than their surface syntax. Matching
$\equiv_{\mathrm{str}}$ to each notation's \emph{native}
equivalence is the reverse-map question of Section~\ref{subsec:reverse-map}, which is partly open.

\begin{corollary}[Internal routing is quotiented]\label{cor:internal-routing-quotient}
Let $m_1$ and $m_2$ differ only in internal places or in silent $\tau$ routing, so that they present the same observable contexts $C(m_1)=C(m_2)$. Then $\Sigma(m_1)=\Sigma(m_2)$.
\end{corollary}

\begin{proof}
Silent transitions contribute no generator, and internal places are XOR mediators that the walk of Definition~\ref{def:contexts-general} passes through. By Theorem~\ref{thm:signature-equivalence}(i) they leave the context structure $C(m)$ unchanged, so Lemma~\ref{lem:sig-complete} gives $\Sigma(m_1)=\Sigma(m_2)$.
\end{proof}

\subsection{Reverse maps and faithfulness}\label{subsec:reverse-map}

The converse direction, recovering a model from a signature, we develop only as far as stating what
each reconstruction requires. The target is a single schema parameterised by the notation
$\mathcal{N}$ and its native behavioural equivalence $\sim_{\mathcal{N}}$. For a forward map
$\Phi_{\mathcal{N}}$ and a proposed reverse map $\Psi_{\mathcal{N}}$, we want $\Phi_{\mathcal{N}}
\circ \Psi_{\mathcal{N}}$ the identity on signatures and $\Psi_{\mathcal{N}} \circ
\Phi_{\mathcal{N}}$ the identity on models up to $\sim_{\mathcal{N}}$. Where each instance is
canonical, and where it weakens to a normal form, is summarised in Table~\ref{tab:faithfulness}.

\begin{table*}[t]
  \centering
  \caption{Reverse maps and the canonicity of each round trip, by notation. SP abbreviates series--parallel throughout.}
  \label{tab:faithfulness}
  \small\begin{tabularx}{\linewidth}{@{}llX@{}}
    \hline
    Notation & Reverse map & Canonicity \\
    \hline
    Causal nets   & read boundaries back as bindings & identity up to isomorphism \\
    Petri nets    & apex hyperedge $+$ glued ports    & up to silent ($\tau$) saturation \\
    Process trees & SP decomposition tree ($\to,+$)   & canonical on $\to,+$; open with XOR/loop \\
    BPMN          & XOR/AND normal form               & normal form, not original gateways \\
    \hline
  \end{tabularx}
\end{table*}

For causal nets the boundary data of the generator cospans are exactly the binding sets
(Section~\ref{sec:causal-nets}), so the round trip is the identity up to isomorphism. For Petri nets
each generator reads back as a transition (apex label, glued ports as its places), the identity up
to $\tau$-saturation. Silent transitions produce no generators, so the reverse map recovers the
$\tau$-free representative. For process trees the $\to,+$ fragment is canonical and constructive
(Corollary~\ref{cor:pt-sp-iff}), while recovering XOR and loop structure from a set of SP posets is
open (Section~\ref{subsec:open-problems}). For BPMN the OR explosion is many-to-one, so a signature
fixes a canonical XOR/AND normal form, not the OR-gateway syntax that may have produced it,
consistent with the denotational reading of Section~\ref{sec:bpmn}.
These four reverse maps are the model-to-model transformations that would slot into
existing modelling toolchains, generalising the BPMN, Petri net, causal net, and process tree
conversions of \citet{kalenkovaProcessMiningUsing2017}.
Synthesis of Petri nets from partial-order languages has been studied in general through
token-flow regions~\cite{bergenthumSynthesisPetriNets2008}. The nets produced by synthesis are
typically hard for an analyst to read~\cite{leemansPartialorderbasedProcessMining2023}, which
motivates recovering the structured notations above instead.

\subsection{Open problems}\label{subsec:open-problems}

\begin{enumerate}
  \item \emph{Process-tree definability.} Characterise which SP-poset languages are of the form
    $\mathrm{sem}_{\mathrm{pt}}(T)$ for some process tree $T$. Process-tree languages are a strict
    subclass of SP-poset languages. The relevant setting is that of series-parallel
    languages~\cite{lodayaSeriesParallelLanguages2000}.
  \item \emph{Canonical process trees.} Even for definable languages, process trees admit no
    canonical form without a normal-form theorem. Define a congruence (building on the equational
    theory of pomsets~\cite{gischerEquationalTheoryPomsets1988}) and prove a minimisation result, in
    particular for the loop operator.
  \item \emph{BPMN faithfulness.} Make precise the sense in which the recovered XOR/AND normal form
    is behaviourally equivalent to the original OR-gateway model.
  \item \emph{POWL as a fifth instance.} The partially ordered workflow language (POWL) composes submodels along arbitrary partial
    orders together with choice and loop operators~\cite{kouraniPOWLPartiallyOrdered2023}. Its
    poset-plus-operators syntax suggests an AND/OR graph in the sense of
    Definition~\ref{def:lm-graph} and hence a signature
    (Definition~\ref{def:generator-cospan-general}). Carrying this out would add a
    partial-order-native notation to the four treated here and relate the certificate of
    Theorem~\ref{thm:signature-equivalence} to POWL's language-preserving workflow-net
    translation~\cite{kouraniTranslatingWorkflowNets2025}.
\end{enumerate}
